
\pdfoutput=1
\documentclass[conference]{IEEEtran}

% Keep this line only when a funding acknowledgment is placed in the
% first-page title footnote.
% \IEEEoverridecommandlockouts

% Official IEEE conference-compatible packages.
\usepackage{cite}
\usepackage{amsmath,amssymb,amsfonts}
\usepackage{graphicx}
\usepackage{textcomp}
\usepackage{xcolor}
\usepackage{booktabs}
\usepackage{multirow}
\usepackage{threeparttable}
\usepackage{adjustbox}
\usepackage{array}
\usepackage{enumitem}
\usepackage{xurl}
\usepackage{dblfloatfix}
\usepackage{capt-of}
\usepackage{hyperref}

\hypersetup{
  hidelinks,
  pdfauthor={Anonymous Authors},
  pdftitle={ToothDINO: A Geometry-Aware Pretrained Vision Encoder for Panoramic Dental Radiographs}
}

% Compact contribution-list spacing within the IEEE page geometry.
\setlist[itemize]{leftmargin=*,topsep=2pt,itemsep=1pt,parsep=0pt,partopsep=0pt}
\emergencystretch=1.5em

\newcommand{\ToothDINO}{\textsc{ToothDINO}}
\newcommand{\DINOB}{\textsc{DINOv3-B}}
\newcommand{\DCT}{\textsc{DINOv3-B-CPT}}
\newcommand{\DCTFull}{\DINOB{} continued pretraining}
\newcommand{\DCTIntro}{\DCTFull{} (\DCT{})}
\newcommand{\DCTControl}{\DCT{} (control)}
\newcommand{\DAVC}{\textsc{DAVC}}
\newcommand{\DXA}{\textsc{DXA}}
\newcommand{\TCC}{\textsc{TCC}}
\newcommand{\ABM}{\textsc{ABM}}
\newcommand{\cmark}{\ensuremath{\checkmark}}
\newcommand{\APb}{\ensuremath{\mathrm{AP}^{b}}}
\newcommand{\APm}{\ensuremath{\mathrm{AP}^{m}}}
\newcommand{\APbFifty}{\ensuremath{\mathrm{AP}^{b}_{50}}}
\newcommand{\APmFifty}{\ensuremath{\mathrm{AP}^{m}_{50}}}
\newcommand{\mIoU}{\ensuremath{\mathrm{mIoU}}}
\newcommand{\mDice}{\ensuremath{\mathrm{mDice}}}
\newcommand{\dataset}[1]{\textsc{#1}}
\newcommand{\best}[1]{\textbf{#1}}
\newcommand{\second}[1]{\underline{#1}}
\newcommand{\bk}[2]{#1~\cite{#2}}
\newcommand{\ourmethod}{\ToothDINO{}-B~}
\newcommand{\CE}{\ensuremath{\operatorname{CE}}}

% Compile-safe figure inclusion. Missing figures are replaced by labeled boxes.
\newcommand{\safeincludegraphics}[2][]{%
  \IfFileExists{#2}{\includegraphics[#1]{#2}}{%
    \begingroup
    \setlength{\fboxsep}{6pt}%
    \fbox{%
      \begin{minipage}[c][0.20\textheight][c]{0.90\linewidth}
        \centering\small Missing figure file:\\[2pt]\detokenize{#2}
      \end{minipage}%
    }%
    \endgroup
  }%
}

\title{ToothDINO: A Geometry-Aware Pretrained Vision Encoder for Panoramic Dental Radiographs}

% Set \anonymousfalse for the camera-ready manuscript.
\newif\ifanonymous
\anonymoustrue

\ifanonymous
\author{\IEEEauthorblockN{Anonymous Authors}}
\else
\author{
\IEEEauthorblockN{Author names omitted for anonymous review}
\IEEEauthorblockA{Affiliations omitted for anonymous review}
}
\fi

\begin{document}
\maketitle

\begin{abstract}
Panoramic dental radiography produces large collections of clinically valuable images, yet expert annotations remain costly and fragmented across classification, detection, and segmentation tasks. Generic self-supervised encoders offer strong transferable representations, but their natural-image view and masking distributions are poorly aligned with the spatially concentrated anatomy of panoramic radiographs. We introduce \ToothDINO{}, a geometry-aware pretrained vision encoder obtained by continuing the official DINOv3 ViT-B/16 checkpoint on 57,232 unlabeled panoramic radiographs. Rather than adding a new backbone or objective, \ToothDINO{} redistributes existing self-supervised supervision through two annotation-free components. Dental-Aware View Construction (\DAVC{}) coordinates radiograph-compatible global views and tooth-centric local views, while Anatomy-Biased Masking (\ABM{}) reallocates masked-token supervision toward anatomy-relevant global patches. Both components operate through sampling and require no tooth boxes, segmentation masks, or task labels. We evaluate the resulting encoder on seven downstream datasets spanning classification, object detection, instance segmentation, and semantic segmentation. Relative to a scale-matched DINOv3-B continued-pretraining control using the same initialization, backbone, corpus, pretraining schedule, and matched rank-$8$ LoRA transfer protocol, \ToothDINO{} improves DENTEX box and mask average precision from 32.7 and 33.9 to 34.1 and 34.7, respectively, and increases MTMD F1 from 81.2 to 86.3. A separate five-seed analysis shows higher mean performance and smaller standard deviations across all reported metrics. The improvements are strongest for anatomically localized targets and more limited for diffuse or heterogeneous targets, suggesting the task-dependent behavior of the fixed anatomical prior.
\end{abstract}

\begin{IEEEkeywords}
panoramic dental radiographs, self-supervised learning, DINOv3, continued pretraining, medical image analysis
\end{IEEEkeywords}


\section{Introduction}
\label{sec:introduction}

Panoramic dental radiographs are routinely acquired for oral screening, diagnosis, and treatment planning because a single examination depicts the teeth, roots, alveolar bone, restorations, lesions, and adjacent maxillofacial structures. The same image can support multiple computational tasks, including image-level diagnosis, tooth and lesion localization, instance delineation, and pixel-level segmentation. However, producing reliable task-specific annotations requires substantial expertise from experienced dental clinicians and is considerably more difficult to scale than routine image acquisition. Consequently, clinical archives contain large collections of panoramic radiographs, while only a relatively small subset is accompanied by annotations suitable for individual downstream tasks. Radiographs acquired using different imaging devices may also exhibit substantial variations in image appearance and quality. The abundance of unlabeled radiographs therefore creates a natural opportunity for self-supervised pretraining, while acquisition heterogeneity further motivates the learning of robust and reusable dental representations \cite{schwendicke2020ai,schneider2023fed_toothseg,rischke2022fldentistry}.

Despite this opportunity, most existing dental image analysis methods are developed as isolated task-specific models, with a separate encoder trained or fine-tuned for each dataset and prediction target. This paradigm repeatedly learns similar anatomical and radiographic patterns from limited task-specific annotations while leaving the visual knowledge contained in large collections of unlabeled panoramic radiographs largely underutilized. A shared pretrained dental encoder provides a more scalable alternative by learning reusable anatomical representations from unlabeled radiographs and subsequently adapting them to classification, detection, and segmentation through task-specific prediction heads.

Self-supervised learning (SSL) has progressed through representative paradigms including contrastive learning, masked image modeling, and teacher--student self-distillation \cite{chen2021empirical,bao2021beit,he2021mae,caron2021emerging}. Contrastive learning emphasizes invariance across image views, masked image modeling captures contextual structure through reconstruction or prediction, and teacher--student self-distillation promotes consistency across transformed views. Among these approaches, Self-Distillation with No Labels (DINO) employs multi-crop teacher--student learning to encourage consistent representations across global and local views, while DINOv2 improves the generality and transferability of these representations through large-scale self-supervised pretraining \cite{caron2021emerging,oquab2023dinov2}. DINOv3 further advances this line of research by providing strong generic representations with spatially detailed dense features, making it a suitable foundation for panoramic dental tasks that require both global image understanding and precise localization \cite{simeoni2025dinov3}. However, a strong generic encoder does not necessarily imply that its original pretraining strategy is well aligned with the geometry and imaging characteristics of panoramic dental radiographs.

Specifically, adapting DINOv3 to panoramic dental radiographs presents two domain-specific challenges: constructing anatomically meaningful global and local views, and allocating masked-token supervision to clinically informative regions. The first challenge concerns view construction and augmentation. Generic random local crops may over-sample background or peripheral regions and therefore provide insufficient coverage of informative tooth-level anatomy. Moreover, color-oriented augmentations developed for natural images are poorly suited to grayscale radiographs, where intensity transformations should instead reflect plausible variations in acquisition conditions and radiographic appearance. The second challenge concerns masked-patch sampling. Clinically informative structures, including teeth, roots, alveolar bone, restorations, and pathology-related regions, are concentrated primarily within the tooth-bearing region, whereas blank margins, background regions, and scanner-dependent artifacts occupy a substantial portion of the image. Uniform mask sampling ignores this nonuniform anatomical distribution and may allocate excessive masked-token supervision to weakly informative regions. Directly continuing DINOv3 pretraining with its original natural-image recipe may therefore fail to make effective use of the anatomical geometry and modality-specific characteristics of panoramic dental radiographs.

To address these challenges, we introduce \ToothDINO{}, a geometry-aware pretrained vision encoder obtained by continuing the official DINOv3 Vision Transformer Base checkpoint on 57,232 unlabeled panoramic dental radiographs. The central idea is geometry-aware allocation of existing self-supervised supervision. Dental-Aware View Construction (\DAVC{}) uses Dental X-ray Augmentation (\DXA{}) and Tooth-Centric Cropping (\TCC{}) to coordinate panorama-level context and tooth-level observations, while Anatomy-Biased Masking (\ABM{}) reallocates masked-token supervision toward anatomy-relevant global patches. Both components are annotation-free sampling procedures that operate before the student--teacher forward pass and require no tooth boxes, segmentation masks, or task labels.

We evaluate the resulting encoder on seven downstream datasets and ten evaluation settings spanning object detection, instance segmentation, semantic segmentation, and single- and multi-label classification. The primary evidence comes from a controlled comparison with a scale-matched DINOv3-B continued-pretraining baseline using the same official initialization, backbone, pretraining corpus, training schedule, downstream heads, and evaluation protocols. Broader supervised and self-supervised encoders are included to assess the competitiveness of the learned dental representation, while differences in model scale and pretraining configuration prevent them from serving as strictly controlled evidence.

Our contributions are threefold:
\begin{itemize}
\item We introduce \ToothDINO{}, a reusable pretrained dental encoder learned from 57,232 unlabeled panoramic radiographs and transferred across classification, object detection, instance segmentation, and semantic segmentation.

\item We formulate dental continued pretraining as geometry-aware supervision allocation. \DAVC{} controls which anatomical content is observed, while \ABM{} controls where masked-token supervision is applied through annotation-free geometric sampling.

\item We conduct controlled and broad evaluations across seven datasets and four task families, together with ablations and five-seed downstream stability analysis, to characterize the gains, task-dependent behavior, and limitations of the proposed anatomical prior.
\end{itemize}


\section{Method}
\label{sec:method}

\subsection{Overview}

\ToothDINO{} adapts the official DINOv3 ViT-B/16 checkpoint to panoramic dental radiographs by aligning the input views and masked-token supervision with panoramic dental geometry. Dental-Aware View Construction (\DAVC{}) controls the global and local image content presented during continued pretraining, while Anatomy-Biased Masking (\ABM{}) increases masked-token supervision in anatomy-relevant regions.

As shown in Fig.~\ref{fig:toothdino-framework}, given an input panoramic radiograph, \DAVC{} first applies a shared medical augmentation and then generates two global views through global cropping and eight tooth-centric local views through Tooth-Centric Cropping (\TCC{}). \ABM{} is applied only to the student global branch. The teacher receives the unmasked global views, whereas the student receives the masked global views together with the tooth-centric local views. After continued pretraining with the DINO, iBOT, and KoLeo objectives used in DINOv3 pretraining, the student backbone is retained as the pretrained dental encoder.

\begin{figure*}[t]
  \centering
  \safeincludegraphics[
    width=\textwidth,
    height=0.34\textheight,
    keepaspectratio
  ]{figures/toothdino_overall_framework666.pdf}
  \caption{\ToothDINO{} continued pretraining. \DAVC{} constructs two global views and eight tooth-centric local views using shared medical augmentation, \TCC{}, and branch-specific \DXA{}. \ABM{} samples anatomy-biased masked positions only for the student global branch.}
  \label{fig:toothdino-framework}
\end{figure*}


\subsection{DINOv3 Continued Pretraining}

\ToothDINO{} follows the original DINOv3 student--teacher framework. The teacher is updated through an exponential moving average of the student parameters. DINO performs cross-view class-token self-distillation, iBOT supervises masked student patch tokens using teacher patch targets, and KoLeo regularizes the representation space \cite{zhou2021ibot,kozachenko1987entropy,simeoni2025dinov3}. The later Gram-anchoring refinement stage is not applied during dental continued pretraining.

Let $V_g$ denote the unmasked global views, $V_\ell$ the tooth-centric local views, and $\widehat{V}_g$ the masked student global views. The teacher and student input sets are:
\begin{equation}
\label{eq:teacher-student-view-sets}
\mathcal{V}_t(x)=V_g,
\qquad
\mathcal{V}_s(x)=\widehat{V}_g\cup V_\ell
\end{equation}

The objective retains the DINO, iBOT, and KoLeo terms used during DINOv3 pretraining:
\begin{equation}
\label{eq:toothdino-objective}
\mathcal{L}_{\mathrm{ToothDINO}}
=
\mathcal{L}_{\mathrm{DINO}}
+
\lambda_{\mathrm{iBOT}}\mathcal{L}_{\mathrm{iBOT}}
+
\lambda_{\mathrm{KoLeo}}\mathcal{L}_{\mathrm{KoLeo}}
\end{equation}

The proposed components change the construction of $V_g$, $V_\ell$, and $\widehat{V}_g$ to redistribute the visual evidence used for self-supervision while leaving these loss terms unchanged.

\DAVC{} and \ABM{} operate as preprocessing components for view construction and mask-position selection. Their additional computation is limited to lightweight augmentation and sampling, and they require no downstream annotations.

\subsection{Dental-Aware View Construction}
\label{sec:davc}

\DAVC{} combines \DXA{} and \TCC{} within a shared anatomical frame. Within \DAVC{}, \DXA{} comprises the shared full-image transformation $T_{\mathrm{med}}$ and the branch-specific crop-level transformations $T_{\mathrm{DXA}}^{(g)}$ and $T_{\mathrm{DXA}}^{(\ell)}$. Given a panoramic radiograph $x$, the shared full-image stage of \DXA{} first produces
\begin{equation}
\label{eq:shared-medical-augmentation}
\tilde{x}=T_{\mathrm{med}}(x)
\end{equation}

Applying this shared transformation before cropping ensures that the global and local branches observe the same transformed dental anatomy.

The two branches are then constructed as:
\begin{align}
v_i^{(g)}
&=
T_{\mathrm{DXA},i}^{(g)}
\left(
\operatorname{Crop}_{g,i}(\tilde{x})
\right),
\quad i\in\{1,2\}
\label{eq:davc-global-views}
\\
v_k^{(\ell)}
&=
T_{\mathrm{DXA},k}^{(\ell)}
\left(
\operatorname{Crop}(\tilde{x};c_k,s_k)
\right),
\quad k\in\{1,\ldots,8\}
\label{eq:davc-local-views}
\end{align}

The resulting view sets are:
\begin{equation}
\label{eq:davc-view-sets}
V_g=\{v_1^{(g)},v_2^{(g)}\},
\qquad
V_\ell=\{v_k^{(\ell)}\}_{k=1}^{8}
\end{equation}

The global views preserve panorama-level tooth-bearing anatomy for stable teacher targets, whereas the local views increase the student's exposure to representative tooth-scale structures.

\subsubsection{Dental X-ray Augmentation}

Dental X-ray Augmentation (\DXA{}) uses radiograph-compatible intensity and geometry transformations without color-specific operations. The shared stage $T_{\mathrm{med}}$ applies gamma $[0.8,1.2]$, contrast $[0.85,1.15]$, brightness $[0.9,1.1]$, and small in-plane rotation $[0,5^\circ]$, while the branch-specific stages adjust the crop-level strength: $T_{\mathrm{DXA}}^{(g)}$ is weaker for panorama-level alignment, and $T_{\mathrm{DXA}}^{(\ell)}$ is stronger for tooth-centric crops through increased local contrast, blur, and Gaussian-noise variation with noise standard deviation in $[0.005,0.02]$. Hue, saturation, and solarization are excluded because they have no meaningful counterpart in grayscale radiography.




\subsubsection{Tooth-Centric Cropping}

\TCC{} determines where the local branch observes the shared augmented panorama. It is a weak dental geometric prior rather than a supervised tooth detector and requires neither tooth boxes nor pixel-level annotations.

Given the shared augmented panorama $\tilde{x}$, \TCC{} first computes
an edge response and then constructs the edge--intensity saliency map:
\begin{align}
E_{\tilde{x}}(u,v)
&=
\left|\nabla_u\tilde{x}(u,v)\right|
+
0.75\left|\nabla_v\tilde{x}(u,v)\right|
\label{eq:tcc-edge-response}
\\
S_{\tilde{x}}(u,v)
&=
E_{\tilde{x}}(u,v)
\left(
0.35
+
0.65\,\psi\bigl(\tilde{x}(u,v)\bigr)
\right)
\pi_v(v)
\label{eq:tcc-saliency}
\end{align}

where $\psi(z)=\operatorname{clip}((z-0.18)/0.52,0,1)$ suppresses near-empty dark regions, and $\pi_v(v)$ is a broad Gaussian vertical prior centered near the expected dental band. Smoothed horizontal and vertical projections of $S_{\tilde{x}}$ estimate the dental extent and its vertical center. In practice, columns whose smoothed energy exceeds a fixed fraction of the maximum define the left and right dental extent, and high-energy rows within this extent define the band center. The estimated band is clipped to a broad central panoramic range to prevent collapse on atypical images. Representative locations are then distributed from left to right along the upper and lower tooth rows, forming the candidate set $\mathcal{C}_{\mathrm{rep}}(\tilde{x})$.

For each local view, the crop center and its bounded spatial jitter are
defined as:
\begin{align}
c_k
&=
c_{q(k)}+\epsilon_k,
\quad
c_{q(k)}\in\mathcal{C}_{\mathrm{rep}}(\tilde{x})
\label{eq:tcc-centers}
\\
\epsilon_k
&\sim
U(-\tau_x,\tau_x)
\times
U(-\tau_y,\tau_y)
\label{eq:tcc-center-jitter}
\end{align}

Bounded spatial jitter maintains crop diversity while reducing background-dominated local observations. In normalized image coordinates, the horizontal energy threshold is $0.48$ of the maximum smoothed column energy, the row threshold is $0.58$ of the maximum smoothed row energy within the estimated extent, and the spatial jitters are set to $\tau_x=\tau_y=0.15$; local crop scales are sampled from $[0.65,1.0]$ after tooth-centric center selection. The two global views use crop scales in $[0.32,1.0]$ and resolution $256\times256$, while the eight tooth-centric local views are represented at $128\times128$.

The fixed \TCC{} and \ABM{} priors are used to preserve stable adaptation from the official DINOv3 initialization; more image-adaptive variants increased training cost and degraded transfer in preliminary trials.







\subsection{Anatomy-Biased Masking}
\label{sec:abm}

\ABM{} receives the two global views $V_g$ generated by \DAVC{} and
constructs an anatomy-biased sampling prior over each global-view patch
grid. It only determines which patch positions are masked in the
student global branch; the teacher receives the corresponding
unmasked global views. Let $(u_a,v_b)$ denote the normalized
coordinates of a patch on the global-view grid. A static dental-band
prior is parameterized by center $(c_x,c_y)$ and spread
$(\sigma_x,\sigma_y)$. The normalized distance, sampling weight, and
probability are:

\begin{align}
d_{ab}^{2}
&=
\frac{(u_a-c_x)^2}{\sigma_x^2}
+
\frac{(v_b-c_y)^2}{\sigma_y^2}
\label{eq:mask-prior-distance}
\\
w_{ab}
&=
1+\beta\exp\left(-\frac{1}{2}d_{ab}^{2}\right)
\label{eq:mask-prior-weight}
\\
p_{ab}
&=
\frac{w_{ab}}{\sum_{m,n}w_{mn}}
\label{eq:mask-prior-probability}
\end{align}

The resulting distribution
$P=\{p_{ab}\}_{a,b}$ corresponds to the anatomy-biased sampling prior
shown in Fig.~\ref{fig:toothdino-framework}. Its uniform component
preserves coverage outside the dental band, whereas the Gaussian
component assigns higher masking probability to patches near the
tooth-bearing region. The prior is not intended to precisely trace the
curved dental arch; instead, it provides a weak elliptical
approximation of the spatial concentration of dental anatomy. Since
global crops are sampled before masking and the deployed prior is
applied in crop coordinates, the fixed center provides a coarse
center-weighted approximation of the typical tooth-bearing distribution
in global views. The prior is fixed and annotation-free and requires no
tooth bounding boxes or segmentation masks.

For a target masking ratio $r\sim U(0.1,0.5)$ on a grid containing
$N$ patches, the number of masked positions is
$M=\lfloor rN\rfloor$. A separate mask is sampled for each global view
from $P$ using weighted sampling without replacement:

\begin{equation}
\label{eq:mask-prior-wor}
\mathcal{M}_i
\sim
\operatorname{WOR}
\left(
M,P
\right),
\qquad
i\in\{1,2\}
\end{equation}

At the sampled positions, the corresponding student patch embeddings
are replaced by the standard mask token, yielding the masked global
views $\widehat{V}_g$. These views are processed by the student encoder
to produce patch-level predictions, whereas the teacher processes the
corresponding unmasked global views $V_g$ to generate patch targets.
This forms the ABM pathway shown in
Fig.~\ref{fig:toothdino-framework}:
$V_g$ is converted into an anatomy-biased sampling prior, the prior
determines $\widehat{V}_g$, and the student predictions at the sampled
positions are supervised through the iBOT loss:
\begin{equation}
\label{eq:ibot-masked-positions}
\mathcal{L}_{\mathrm{iBOT}}
=
\frac{1}{\sum_i|\mathcal{M}_i|}
\sum_{i=1}^{2}
\sum_{(a,b)\in\mathcal{M}_i}
H\left(q^t_{i,ab},p^s_{i,ab}\right)
\end{equation}
where $q^t_{i,ab}$ is the teacher patch target from the $i$th unmasked
global view and $p^s_{i,ab}$ is the corresponding student prediction
from the masked global view. We use $\beta=2.5$,
$(c_x,c_y)=(0.5,0.5)$, and
$(\sigma_x,\sigma_y)=(0.35,0.20)$. Thus, \ABM{} reallocates iBOT
supervision toward patches near the tooth-bearing region while
preserving nonzero sampling probability throughout the global view.






\section{Experiments}
\label{sec:experiments}

\subsection{Datasets}
\label{sec:datasets}

The continued-pretraining corpus contains 57,232 panoramic dental radiographs collected from the training sets of public datasets and a de-identified institutional clinical collection. Combining public and institutional data increases the diversity of image sources and acquisition conditions represented during continued pretraining. Only training images from the corresponding public datasets are included in continued pretraining. Their annotations are not used, and no validation or test image from any downstream benchmark is included in the pretraining corpus. The institutional radiographs were de-identified before analysis. % TODO: Add institutional retrospective-collection, ethics-approval, and informed-consent-waiver statement if approved for anonymous review.

Downstream transfer is evaluated on seven datasets covering four task families: \dataset{DENTEX}~\cite{hamamci2023dentex}, \dataset{ADCD}~\cite{muresanu2024adcd}, \dataset{CariesXrays}~\cite{chen2024cariesxrays}, and \dataset{OralXrays-9}~\cite{chen2025OralXrays} for object detection and instance segmentation; a multi-center panoramic collection~\cite{li2024multicenter} and \dataset{DC1000}~\cite{wang2023multi} for semantic segmentation; and \dataset{MTMD}~\cite{mtm_mendeley}, \dataset{OralXrays-9}~\cite{chen2025OralXrays}, and \dataset{ADCD}~\cite{muresanu2024adcd} for single- or multi-label classification. Several datasets support more than one prediction formulation, resulting in ten evaluation settings in total.


\textbf{Multi-label Construction.}
For \dataset{OralXrays-9} and \dataset{ADCD}, we derive image-level
multi-label targets from their COCO-format instance annotations. For
each image, the target label set is constructed by collecting the
categories of all valid annotated instances:
\begin{equation}
Y_i=
\operatorname{unique}
\left\{
\phi(a_{\mathrm{cat}})
\mid
a_{\mathrm{img}}=i
\right\}
\label{eq:multilabel-construction}
\end{equation}
where $\phi(\cdot)$ maps the original category IDs to contiguous class
indices. Multiple instances belonging to the same category contribute
a single image-level label, whereas different categories are retained
jointly. Images without valid instance annotations are excluded. The
conversion is performed separately within the original training,
validation, and test splits to avoid cross-split label leakage.



\subsection{Experimental Setup}
\label{sec:experimental-setup}

We initialize from official DINOv3 ViT-B/16 weights \cite{simeoni2025dinov3} and train for 200 epochs with AdamW, a base learning rate of $10^{-4}$, and batch size $64$ per graphics processing unit (GPU). Training uses weight decay $0.04$ to $0.4$, a cosine learning-rate schedule with 10 warm-up epochs, $\lambda_{\mathrm{iBOT}}=1.0$, and $\lambda_{\mathrm{KoLeo}}=0.1$. Following the pipeline in Fig.~\ref{fig:toothdino-framework}, each panorama yields two global views at $256\times256$ through global cropping and \DXA{}, together with eight tooth-centric local views at $128\times128$ through \TCC{} and local \DXA{}. The teacher receives the unmasked global views; the student receives the anatomy-masked global views and all tooth-centric local views.


Within each benchmark, all encoders follow the same evaluation
protocol, and the main comparison tables use a common fixed downstream
random seed. Classification uses MMPreTrain \cite{2023mmpretrain} with
a linear classification head, $512\times512$ inputs, batch size $16$,
and 50 epochs. Detection and instance segmentation use MMDetection
\cite{mmdetection} with a Mask Region-Based Convolutional Neural Network
(Mask R-CNN) and Feature Pyramid Network (FPN)
\cite{he2017maskrcnn,lin2017fpn}, $1333\times800$ inputs, batch size
$4$, and a 12-epoch schedule. Semantic segmentation uses MMSegmentation
\cite{mmseg2020} with an FPN and Unified Perceptual Parsing Network
(UPerNet)-style decoder together with an auxiliary fully convolutional
network (FCN) head \cite{lin2017fpn,xiao2018upernet},
$1024\times1024$ crops, batch size $4$, and 20k iterations. The
controlled \DCT{}--\ToothDINO{} comparison uses matched rank-$8$ LoRA
adaptation, the same optimizer, and the same task-specific
learning-rate schedule. For the stability analysis, this controlled
comparison is repeated over five downstream random seeds, with results
reported as mean $\pm$ standard deviation in
Table~\ref{tab:seed-stability}. The broader supervised and
self-supervised baselines use full encoder fine-tuning. Names of public
encoders and datasets follow their source publications; \DCT{} denotes
our scale-matched continued-pretraining control initialized from the
official DINOv3-B weights.
LoRA is applied to the attention and feed-forward projections
(\mbox{qkv}, \mbox{proj}, \mbox{fc1}, and \mbox{fc2}) with
rank $r=8$, scaling factor $\alpha=16$, and dropout $0.05$. The
remaining encoder parameters are frozen, while the task-specific
prediction head, neck, or decoder is fully optimized.


Detection and instance segmentation are evaluated using Common Objects in Context (COCO) box AP (\APb{}), box AP$_{50}$ (\APbFifty{}), mask AP (\APm{}), and mask AP$_{50}$ (\APmFifty{}) \cite{lin2014coco}. Semantic segmentation is evaluated using mean intersection over union (\mIoU{}) and mean Dice score (\mDice{}). Segmentation metrics are averaged over all configured classes, including the background class. For the single-label \dataset{MTMD} task, we report accuracy and macro-averaged F1. For the multi-label \dataset{OralXrays-9} and \dataset{ADCD} tasks, precision, recall, and F1 are macro-averaged across classes. Their mean average precision is computed by averaging class-wise AP, and class probabilities are binarized using a fixed threshold of $0.5$:
\begin{equation}
\label{eq:classification-metrics}
\mathrm{mAP}
=
\frac{1}{C}\sum_{c=1}^{C}\mathrm{AP}_{c},
\qquad
\hat{y}_{i,c}
=
\mathbb{I}\!\left[p_{i,c}\geq0.5\right]
\end{equation}

\subsection{Main Results}
\label{sec:main-results}
Unless otherwise noted, bold and underlined entries denote the best and second-best results within each comparison, respectively.

We evaluate \ToothDINO{} from two complementary perspectives. The
matched comparison with \DCT{} isolates the contribution of the
proposed geometry-aware supervision allocation because the two models
share the same initialization, backbone, pretraining corpus, training
schedule, downstream heads, LoRA configuration, and evaluation
protocols. The remaining supervised and self-supervised encoders are
included as external reference comparisons to contextualize the
overall competitiveness of the learned representation.

\subsubsection{Detection and Instance Segmentation}
Tables~\ref{tab:detection-main-a} and~\ref{tab:detection-main-b} summarize the detection and instance-segmentation results on four benchmarks. On \dataset{DENTEX}, \ourmethod improves the \DCT{} control in both box and mask AP, increasing them from $32.7$ and $33.9$ to $34.1$ and $34.7$, respectively. On \dataset{ADCD}, it also improves the matched control across localization metrics. These gains indicate that geometry-aware supervision allocation is most helpful when the downstream target is spatially concentrated within the tooth-bearing region.

\begin{table}[t]
  \centering
  \scriptsize
  \setlength{\tabcolsep}{1.25pt}
  \renewcommand{\arraystretch}{1.03}
  \caption{Detection and instance segmentation on DENTEX and ADCD.}
  \label{tab:detection-main-a}
  \begin{adjustbox}{max width=\columnwidth}
  \begin{tabular}{lcccccccc}
    \toprule
    \multirow{2}{*}{Encoder}
    & \multicolumn{4}{c}{DENTEX}
    & \multicolumn{4}{c}{ADCD} \\
    \cmidrule(lr){2-5}\cmidrule(lr){6-9}
    & \APb & \APbFifty{} & \APm & \APmFifty{}
    & \APb & \APbFifty{} & \APm & \APmFifty{} \\
    \midrule
    \multicolumn{9}{c}{\textbf{Supervised pretraining}} \\
    \midrule
    \bk{ResNet-50}{he2016resnet}
    & 26.4 & 42.4 & 26.5 & 42.7
    & 19.6 & \second{46.5} & 19.6 & 46.6 \\
    \bk{ConvNeXt-B}{liu2022convnext}
    & 30.6 & 50.2 & 31.6 & 50.3
    & 19.4 & 46.2 & 19.2 & 45.8 \\
    \bk{Swin-B}{liu2021swin}
    & 32.0 & 52.8 & 32.5 & 52.4
    & 19.4 & 46.3 & 18.9 & 45.2 \\
    \bk{DeiT-B}{touvron2021training}
    & 21.6 & 37.7 & 22.0 & 37.4
    & 17.2 & 43.1 & 17.1 & 43.1 \\
    \bk{AugReg-L}{steiner2021how}
    & 21.8 & 38.2 & 22.0 & 38.3
    & 17.1 & 42.2 & 16.6 & 41.5 \\
    \midrule
    \multicolumn{9}{c}{\textbf{Self-supervised pretraining}} \\
    \midrule
    \bk{MAE-B}{he2021mae}
    & 22.0 & 36.6 & 21.9 & 36.6
    & 17.3 & 42.3 & 17.9 & 43.5 \\
    \bk{MoCov3}{chen2021empirical}
    & 28.9 & 48.2 & 29.9 & 46.5
    & 19.3 & 45.4 & 19.3 & 45.9 \\
    \bk{BiomedCLIP}{zhang2023biomedclip}
    & 23.1 & 39.9 & 23.8 & 40.0
    & 17.3 & 43.1 & 17.4 & 42.9 \\
    \bk{UniPerceiver-B}{zhu2022uniperceiver}
    & 24.9 & 43.3 & 24.8 & 43.2
    & 18.3 & 44.6 & 18.3 & 45.4 \\
    \bk{BEiT-B}{bao2021beit}
    & 19.2 & 35.3 & 18.1 & 35.6
    & 15.3 & 37.4 & 15.1 & 37.7 \\
    \bk{DINOv2-S}{oquab2023dinov2}
    & 27.3 & 47.7 & 28.1 & 46.5
    & 18.5 & 45.0 & 18.7 & 45.1 \\
    \bk{DINOv2-B}{oquab2023dinov2}
    & 28.7 & 48.7 & 29.2 & 48.5
    & 18.8 & 44.7 & 18.4 & 44.5 \\
    \bk{DINOv2-L}{oquab2023dinov2}
    & 32.1 & 53.7 & 34.1 & 53.3
    & 18.6 & 45.4 & 18.4 & 45.3 \\
    \bk{DINOv3-S}{simeoni2025dinov3}
    & 30.6 & 48.8 & 31.3 & 48.6
    & 18.7 & 44.5 & 18.9 & 44.9 \\
    \bk{\DINOB{}}{simeoni2025dinov3}
    & 31.8 & 51.4 & 32.6 & 50.8
    & 19.2 & 45.1 & 19.4 & 45.2 \\
    \bk{DINOv3-L}{simeoni2025dinov3}
    & \best{34.4} & \second{54.7} & \best{35.6} & \best{55.3}
    & 19.0 & 44.3 & 18.8 & 43.6 \\
    \DCTControl
    & 32.7 & 53.9 & 33.9 & 53.2
    & \second{19.8} & 46.1 & \second{19.9} & \second{46.8} \\
    \addlinespace[2pt]
    \ourmethod (ours)
    & \second{34.1} & \best{55.6} & \second{34.7} & \second{55.2}
    & \best{20.0} & \best{47.2} & \best{20.4} & \best{47.7} \\
    \bottomrule
  \end{tabular}
  \end{adjustbox}
\end{table}

\begin{table}[t]
  \centering
  \scriptsize
  \setlength{\tabcolsep}{1.25pt}
  \renewcommand{\arraystretch}{1.03}
  \caption{Detection and instance segmentation on CariesXrays and OralXrays-9.}
  \label{tab:detection-main-b}
  \begin{adjustbox}{max width=\columnwidth}
  \begin{tabular}{lcccccccc}
    \toprule
    \multirow{2}{*}{Encoder}
    & \multicolumn{4}{c}{CariesXrays}
    & \multicolumn{4}{c}{OralXrays-9} \\
    \cmidrule(lr){2-5}\cmidrule(lr){6-9}
    & \APb & \APbFifty{} & \APm & \APmFifty{}
    & \APb & \APbFifty{} & \APm & \APmFifty{} \\
    \midrule
    \multicolumn{9}{c}{\textbf{Supervised pretraining}} \\
    \midrule
    \bk{ResNet-50}{he2016resnet}
    & \best{46.8} & \best{82.2} & \best{21.8} & \best{40.4}
    & \best{68.4} & \best{91.3} & \best{68.0} & \best{91.4} \\
    \bk{ConvNeXt-B}{liu2022convnext}
    & \second{46.0} & 80.4 & \second{21.1} & 39.3
    & 67.2 & 90.1 & 67.1 & 90.1 \\
    \bk{Swin-B}{liu2021swin}
    & 43.0 & 77.8 & 19.9 & 38.3
    & 66.4 & 90.3 & 66.1 & 90.3 \\
    \bk{DeiT-B}{touvron2021training}
    & 41.3 & 77.2 & 15.1 & 31.7
    & 67.4 & \second{90.8} & 67.2 & 90.7 \\
    \bk{AugReg-L}{steiner2021how}
    & 39.8 & 76.3 & 12.6 & 28.1
    & 66.6 & 90.4 & 66.4 & 90.5 \\
    \midrule
    \multicolumn{9}{c}{\textbf{Self-supervised pretraining}} \\
    \midrule
    \bk{MAE-B}{he2021mae}
    & 42.8 & 78.7 & 20.0 & 38.7
    & 65.3 & 90.5 & 64.9 & 90.6 \\
    \bk{MoCov3}{chen2021empirical}
    & 44.9 & \second{81.8} & 20.6 & \second{40.0}
    & \second{67.6} & \second{90.8} & 67.2 & \second{91.0} \\
    \bk{BiomedCLIP}{zhang2023biomedclip}
    & 37.0 & 73.0 & 17.1 & 35.1
    & 63.6 & 89.1 & 63.0 & 89.2 \\
    \bk{UniPerceiver-B}{zhu2022uniperceiver}
    & 42.2 & 78.6 & 14.9 & 31.3
    & \second{67.6} & 90.4 & \second{67.5} & 90.6 \\
    \bk{BEiT-B}{bao2021beit}
    & 40.2 & 76.6 & 13.5 & 28.8
    & 66.3 & 90.1 & 66.1 & 90.3 \\
    \bk{DINOv2-S}{oquab2023dinov2}
    & 39.4 & 75.4 & 18.5 & 37.2
    & 65.6 & 89.5 & 65.5 & 89.5 \\
    \bk{DINOv2-B}{oquab2023dinov2}
    & 40.8 & 76.9 & 19.0 & 37.8
    & 65.7 & 89.2 & 65.6 & 89.4 \\
    \bk{DINOv2-L}{oquab2023dinov2}
    & 42.1 & 77.6 & 19.7 & 38.3
    & 66.0 & 89.3 & 65.8 & 89.5 \\
    \bk{DINOv3-S}{simeoni2025dinov3}
    & 40.8 & 77.0 & 18.8 & 37.8
    & 66.1 & 90.3 & 66.2 & 90.4 \\
    \bk{\DINOB{}}{simeoni2025dinov3}
    & 41.9 & 79.0 & 19.7 & 38.8
    & 65.9 & 89.8 & 65.7 & 89.8 \\
    \bk{DINOv3-L}{simeoni2025dinov3}
    & 43.3 & 78.9 & 20.4 & 38.7
    & 65.8 & 89.3 & 65.6 & 89.4 \\
    \DCTControl
    & 43.3 & 80.2 & 19.3 & 38.5
    & 67.1 & 90.7 & 66.9 & \second{91.0} \\
    \addlinespace[2pt]
    \ourmethod (ours)
    & 44.3 & 81.0 & 20.5 & 39.3
    & 67.0 & 90.7 & 67.0 & 90.8 \\
    \bottomrule
  \end{tabular}
  \end{adjustbox}
\end{table}

On \dataset{CariesXrays}, \ourmethod improves box AP by $1.0$ point and mask AP by $1.2$ points over \DCT{}. On \dataset{OralXrays-9}, the aggregate box and mask AP values remain nearly unchanged. These results show that the effect of geometry-aware supervision is dataset dependent rather than uniformly positive.

\subsubsection{Semantic Segmentation}
Table~\ref{tab:segmentation-main} reports the semantic-segmentation results. On \dataset{DC1000}, \ourmethod improves the \DCT{} control from $77.4/85.5$ to $78.0/86.0$ in \mIoU{}/\mDice{}. Impacted-tooth and periodontitis segmentation remain comparable to the control, indicating that the benefit does not extend uniformly across all dense-prediction targets.

\begin{table}[t]
  \centering
  \scriptsize
  \setlength{\tabcolsep}{2.2pt}
  \renewcommand{\arraystretch}{1.05}
  \caption{Semantic segmentation on panoramic dental datasets.}
  \label{tab:segmentation-main}
  \begin{adjustbox}{max width=\columnwidth}
  \begin{tabular}{lcccccc}
    \toprule
    \multirow{2}{*}{Encoder}
    & \multicolumn{2}{c}{Impacted}
    & \multicolumn{2}{c}{Periodontitis}
    & \multicolumn{2}{c}{DC1000} \\
    \cmidrule(lr){2-3}
    \cmidrule(lr){4-5}
    \cmidrule(lr){6-7}
    & \mIoU & \mDice
    & \mIoU & \mDice
    & \mIoU & \mDice \\
    \midrule

    \multicolumn{7}{c}{\textbf{Supervised pretraining}} \\
    \midrule
    \bk{ViT-B}{dosovitskiy2020image}
      & 88.0 & 93.2 & 88.0 & 93.2 & 76.9 & 85.1 \\
    \bk{Twins-S}{chu2021twins}
      & 87.1 & 92.6 & 94.7 & 97.2 & 78.0 & 85.9 \\
    \bk{Swin-B}{liu2021swin}
      & 81.6 & 88.8 & 95.9 & 97.9 & 70.0 & 78.8 \\
    \bk{ResNeXt-101}{xie2017resnext}
      & 89.4 & \second{94.1} & 95.7 & 97.8 & 77.7 & 85.7 \\
    \bk{ResNet-50}{he2016resnet}
      & 86.5 & 92.2 & 94.9 & 97.3 & 75.2 & 83.7 \\
    \bk{ResNet-101}{he2016resnet}
      & \second{89.7} & \best{94.3} & 95.7 & 97.8 & 77.9 & 85.9 \\
    \bk{ResNeSt-101}{zhang2022resnest}
      & \best{89.8} & \best{94.3} & 95.5 & 97.7 & 77.8 & 85.8 \\
    \bk{PoolFormer-S12}{yu2022poolformer}
      & \best{89.8} & \best{94.3}
      & \best{96.3} & \best{98.1}
      & 77.4 & 85.5 \\
    \bk{MSCAN-T}{guo2022segnext}
      & 87.3 & 92.8 & 94.8 & 97.3 & 76.6 & 84.8 \\
    \bk{HRNet-W48}{wang2020deep}
      & 87.9 & 93.1 & 94.9 & 97.3 & 78.0 & 86.0 \\
    \bk{ConvNeXt-B}{liu2022convnext}
      & \best{89.8} & \best{94.3}
      & \best{96.3} & \best{98.1}
      & \second{79.2} & \second{86.9} \\

    \midrule
    \multicolumn{7}{c}{\textbf{Self-supervised pretraining}} \\
    \midrule
    \bk{MAE-B}{he2021mae}
      & 87.4 & 92.8 & 87.6 & 93.0 & 75.2 & 83.7 \\
    \bk{BEiT-B}{bao2021beit}
      & 88.8 & 93.7 & 88.6 & 93.6 & 78.6 & 86.5 \\
    \bk{DINOv2-S}{oquab2023dinov2}
      & 87.4 & 92.8 & 96.1 & \second{98.0} & 78.4 & 86.3 \\
    \bk{DINOv2-B}{oquab2023dinov2}
      & 87.7 & 93.0 & 96.1 & \second{98.0}
      & \best{79.6} & \best{87.2} \\
    \bk{DINOv2-L}{oquab2023dinov2}
      & 87.4 & 92.8 & 95.7 & 97.8 & 77.9 & 85.9 \\
    \bk{DINOv3-S}{simeoni2025dinov3}
      & 87.9 & 93.1 & 96.0 & 97.9 & 76.6 & 84.8 \\
    \bk{\DINOB{}}{simeoni2025dinov3}
      & 88.7 & 93.6 & \second{96.2} & \second{98.0}
      & 77.4 & 85.5 \\
    \bk{DINOv3-L}{simeoni2025dinov3}
      & 89.0 & 93.9 & \best{96.3} & \best{98.1}
      & 78.7 & 86.5 \\
    \DCTControl
      & 88.8 & 93.7 & 96.0 & 97.9 & 77.4 & 85.5 \\

    \addlinespace[2pt]
    \ourmethod{} (ours)
      & 88.7 & 93.6 & 96.1 & \second{98.0}
      & 78.0 & 86.0 \\
    \bottomrule
  \end{tabular}
  \end{adjustbox}
\end{table}


\subsubsection{Classification}
Table~\ref{tab:classification-main} reports the classification results. On \dataset{MTMD}, \ourmethod improves accuracy from $91.3$ to $93.6$ and F1 score from $81.2$ to $86.3$ relative to \DCT{}. On \dataset{OralXrays-9}, it improves F1 while remaining slightly below \DCT{} in mAP. On \dataset{ADCD}, the continued-pretraining control remains stronger, further showing that the proposed spatial prior is not uniformly beneficial across heterogeneous image-level targets.

\begin{table}[t]
  \centering
  \scriptsize
  \setlength{\tabcolsep}{2.0pt}
  \renewcommand{\arraystretch}{1.05}
  \caption{Classification on panoramic dental datasets.}
  \label{tab:classification-main}
  \begin{threeparttable}
    \begin{adjustbox}{max width=\columnwidth}
    \begin{tabular}{lcccccc}
      \toprule
      \multirow{2}{*}{Encoder}
      & \multicolumn{2}{c}{MTMD}
      & \multicolumn{2}{c}{OralXrays-9}
      & \multicolumn{2}{c}{ADCD} \\
      \cmidrule(lr){2-3}
      \cmidrule(lr){4-5}
      \cmidrule(lr){6-7}
      & Acc & F1
      & mAP & F1
      & mAP & F1 \\
      \midrule

      \multicolumn{7}{c}{\textbf{Supervised pretraining}} \\
      \midrule
      \bk{ResNet-50}{he2016resnet}
        & 90.9 & 81.0 & 95.0 & 84.6 & 68.7 & 61.9 \\
      \bk{Swin-B}{liu2021swin}
        & 90.9 & 81.7 & 95.0 & 84.1 & 68.5 & 61.1 \\
      \bk{SwinV2-B}{liu2022swinv2}
        & 84.5 & 57.8 & 92.0 & 79.5 & 64.4 & 59.4 \\
      \bk{ConvNeXt-B}{liu2022convnext}
        & 89.0 & 74.2 & 90.7 & 80.3 & 67.0 & 59.5 \\
      \bk{EfficientNet-B4}{tan2019efficientnet}
        & 76.1 & 39.1 & 92.0 & 80.4 & 66.6 & 58.0 \\
      \bk{DeiT-B}{touvron2021training}
        & 79.5 & 46.3 & 84.8 & 75.9 & 63.5 & 55.2 \\
      \bk{DeiT3-B}{touvron2022deit}
        & 79.2 & 43.2 & 86.8 & 79.5 & 67.6 & 59.5 \\
      \bk{ViT-B}{dosovitskiy2020image}
        & 77.7 & 48.7 & 84.1 & 75.9 & 64.4 & 48.9 \\
      \bk{RegNetX-8GF}{radosavovic2020designing}
        & 83.0 & 54.5 & 94.1 & 83.7 & 68.7 & 61.4 \\

      \midrule
      \multicolumn{7}{c}{\textbf{Self-supervised pretraining}} \\
      \midrule
      \bk{BEiT-B}{bao2021beit}
        & 82.2 & 52.5 & 93.7 & 85.1 & 65.8 & 60.6 \\
      \bk{DINOv2-S}{oquab2023dinov2}
        & \second{92.0} & 82.7
        & 95.0 & 83.7
        & 71.2 & 62.7 \\
      \bk{DINOv2-B}{oquab2023dinov2}
        & 87.5 & 71.8
        & 95.3 & 84.6
        & 71.0 & 62.5 \\
      \bk{DINOv2-L}{oquab2023dinov2}
        & 86.0 & 62.6
        & 95.7 & 85.8
        & 72.5 & 63.6 \\
      \bk{DINOv3-S}{simeoni2025dinov3}
        & 88.3 & 77.5
        & 95.1 & \second{86.0}
        & 73.1 & 63.7 \\
      \bk{\DINOB{}}{simeoni2025dinov3}
        & 89.0 & 73.0
        & 95.5 & 85.1
        & 70.3 & \second{64.4} \\
      \bk{DINOv3-L}{simeoni2025dinov3}
        & 91.7 & \second{85.4}
        & 95.5 & 85.0
        & 72.5 & 64.2 \\
      \DCTControl
        & 91.3 & 81.2
        & \best{96.0} & 85.2
        & \best{75.2} & \best{68.5} \\

      \addlinespace[2pt]
      \ourmethod{} (ours)
        & \best{93.6} & \best{86.3}
        & \second{95.8} & \best{86.3}
        & \second{74.1} & 63.5 \\
      \bottomrule
    \end{tabular}
    \end{adjustbox}

  %   \begin{tablenotes}[flushleft]
  %     \scriptsize
  %     \item[a] MTMD is a single-label classification dataset; therefore,
  %     accuracy is used as its primary metric instead of mAP.
  %   \end{tablenotes}
  \end{threeparttable}
\end{table}



\subsection{Stability across Downstream Seeds}
\label{sec:seed-stability}

We repeat the controlled \DCT{}--\ToothDINO{} comparison over five
downstream random seeds using the same rank-$8$ LoRA adaptation,
data splits, task heads, and training schedules. As shown in
Table~\ref{tab:seed-stability}, \ToothDINO{} achieves higher mean
performance than \DCT{} on all reported metrics while exhibiting
lower variability across seeds. The improvements are observed across
detection, instance segmentation, semantic segmentation, and
classification, indicating that they are not driven by a single
favorable random seed.

\begin{table}[t]
  \centering
  \scriptsize
  \setlength{\tabcolsep}{4.2pt}
  \renewcommand{\arraystretch}{1.08}
  \caption{Five-seed downstream stability under matched rank-$8$ LoRA,
  reported as mean $\pm$ standard deviation.}
  \label{tab:seed-stability}
  \begin{adjustbox}{max width=\columnwidth}
  \begin{tabular}{llcc}
    \toprule
    \textbf{Task}
    & \textbf{Metric}
    & \textbf{\DCT{}}
    & \textbf{\ToothDINO{}} \\
    \midrule
    \multirow{2}{*}{\dataset{DENTEX} detection}
    & \APb{}
    & $32.84 \pm 0.45$
    & $\mathbf{34.02 \pm 0.30}$ \\
    & \APm{}
    & $33.67 \pm 0.50$
    & $\mathbf{34.86 \pm 0.40}$ \\
    \addlinespace[2pt]

    \multirow{2}{*}{\dataset{CariesXrays} detection}
    & \APb{}
    & $43.58 \pm 0.55$
    & $\mathbf{44.18 \pm 0.30}$ \\
    & \APm{}
    & $19.16 \pm 0.35$
    & $\mathbf{20.57 \pm 0.20}$ \\
    \addlinespace[2pt]

    \multirow{2}{*}{\dataset{DC1000} segmentation}
    & \mIoU{}
    & $77.48 \pm 0.32$
    & $\mathbf{77.94 \pm 0.29}$ \\
    & \mDice{}
    & $85.39 \pm 0.27$
    & $\mathbf{86.10 \pm 0.24}$ \\
    \addlinespace[2pt]

    \multirow{2}{*}{\dataset{MTMD} classification}
    & Acc.
    & $91.55 \pm 0.65$
    & $\mathbf{93.34 \pm 0.58}$ \\
    & F1
    & $81.76 \pm 1.15$
    & $\mathbf{85.88 \pm 0.95}$ \\
    \bottomrule
  \end{tabular}
  \end{adjustbox}
\end{table}



\subsection{Ablation Study}
\label{sec:ablation-study}
We evaluate the two proposed components using four configurations:
the continued-pretraining control, \ABM{} alone, \DAVC{} alone, and
their combination. This design isolates the individual contribution
of each component and also evaluates whether their effects are
complementary. \DAVC{} is treated as an integrated view-construction
component comprising shared medical augmentation, branch-specific
\DXA{}, and \TCC{}.

Box-detection and instance-segmentation ablations are reported in Table~\ref{tab:detection-ablation}. The factorial comparison shows that \ABM{} is effective even without \DAVC{}, especially on \dataset{DENTEX}, where it improves all four metrics over the control. \DAVC{} alone is more effective on ADCD box detection and CariesXrays mask prediction, whereas the complete model obtains the strongest DENTEX results and several complementary gains across the remaining localization benchmarks.

These results indicate that \DAVC{} and \ABM{} provide distinct but partly complementary forms of supervision allocation: \DAVC{} changes which regions are observed, whereas \ABM{} changes which global patches receive masked-token supervision.

\begin{table}[!t]
  \centering
  \scriptsize
  \setlength{\tabcolsep}{2.7pt}
  \renewcommand{\arraystretch}{1.04}
  \caption{Ablation of \DAVC{} and \ABM{} for detection and instance segmentation.}
  \label{tab:detection-ablation}
  \begin{adjustbox}{max width=\columnwidth}
  \begin{tabular}{lcc*{4}{c}}
    \toprule
    \textbf{Dataset} & \DAVC{} & \ABM{}
    & \APb{} & \APbFifty{} & \APm{} & \APmFifty{} \\
    \midrule
    \multirow{4}{*}{ADCD}
    &  &  & 19.8 & 46.1 & 19.9 & 46.8 \\
    &  & \cmark & 19.6 & \second{47.5} & 19.8 & 47.6 \\
    & \cmark &  & \best{20.3} & \best{48.1} & \second{20.2} & \best{48.6} \\
    & \cmark & \cmark & \second{20.0} & 47.2 & \best{20.4} & \second{47.7} \\
    \midrule
    \multirow{4}{*}{CariesXrays}
    &  &  & 43.3 & 80.2 & 19.3 & 38.5 \\
    &  & \cmark & 43.9 & 80.4 & 20.0 & 38.9 \\
    & \cmark &  & \second{44.2} & \second{80.5} & \best{20.7} & \best{39.9} \\
    & \cmark & \cmark & \best{44.3} & \best{81.0} & \second{20.5} & \second{39.3} \\
    \midrule
    \multirow{4}{*}{OralXrays-9}
    &  &  & \best{67.1} & \second{90.7} & \second{66.9} & \best{91.0} \\
    &  & \cmark & \second{67.0} & \best{90.8} & 66.6 & \best{91.0} \\
    & \cmark &  & \second{67.0} & \second{90.7} & 66.5 & \second{90.9} \\
    & \cmark & \cmark & \second{67.0} & \second{90.7} & \best{67.0} & 90.8 \\
    \midrule
    \multirow{4}{*}{DENTEX}
    &  &  & 32.7 & 53.9 & 33.9 & 53.2 \\
    &  & \cmark & \second{33.5} & \second{54.8} & \second{34.0} & \second{54.2} \\
    & \cmark &  & 32.8 & 53.8 & 33.3 & 52.9 \\
    & \cmark & \cmark & \best{34.1} & \best{55.6} & \best{34.7} & \best{55.2} \\
    \bottomrule
  \end{tabular}
  \end{adjustbox}
\end{table}

The semantic-segmentation ablation in Table~\ref{tab:segmentation-ablation} further demonstrates task dependence. \ABM{} alone is strongest for periodontitis, \DAVC{} and the full model are tied on \dataset{DC1000}, and the control remains strongest for impacted-tooth segmentation. Thus, the preferred component allocation varies with the spatial structure of the target.

\begin{table}[t]
  \centering
  \scriptsize
  \setlength{\tabcolsep}{3.0pt}
  \caption{Ablation of \DAVC{} and \ABM{} for semantic segmentation.}
  \label{tab:segmentation-ablation}
  \begin{adjustbox}{max width=\columnwidth}
  \begin{tabular}{cc*{6}{c}}
    \toprule
    \multicolumn{2}{c}{\textbf{Components}}
    & \multicolumn{2}{c}{Impacted}
    & \multicolumn{2}{c}{Periodontitis}
    & \multicolumn{2}{c}{DC1000} \\
    \cmidrule(lr){1-2}
    \cmidrule(lr){3-4}
    \cmidrule(lr){5-6}
    \cmidrule(lr){7-8}
    \DAVC{} & \ABM{} & \mIoU & \mDice & \mIoU & \mDice & \mIoU & \mDice \\
    \midrule
     &  & \best{88.8} & \best{93.7} & 96.0 & \second{97.9} & 77.4 & 85.5 \\
     & \cmark & 88.2 & 93.3 & \best{96.2} & \best{98.0} & \second{77.9} & \second{85.9} \\
    \cmark &  & 88.5 & 93.5 & 96.0 & \second{97.9} & \best{78.0} & \best{86.0} \\
    \cmark & \cmark & \second{88.7} & \second{93.6} & \second{96.1} & \best{98.0} & \best{78.0} & \best{86.0} \\
    \bottomrule
  \end{tabular}
  \end{adjustbox}
\end{table}

The classification ablation in Table~\ref{tab:classification-ablation} shows a different pattern. \ABM{} alone achieves the highest \dataset{MTMD} accuracy and F1, the full model achieves the highest \dataset{OralXrays-9} F1, and the \DCT{} control remains strongest on \dataset{ADCD}. These results show that isolated \ABM{} is useful, but the best configuration depends on the downstream label structure.

\begin{table}[t]
  \centering
  \scriptsize
  \setlength{\tabcolsep}{3.0pt}
  \caption{Ablation of \DAVC{} and \ABM{} for classification.}
  \label{tab:classification-ablation}
  \begin{threeparttable}
  \begin{adjustbox}{max width=\columnwidth}
  \begin{tabular}{cc*{6}{c}}
    \toprule
    \multicolumn{2}{c}{\textbf{Components}}
    & \multicolumn{2}{c}{OralXrays-9}
    & \multicolumn{2}{c}{ADCD}
    & \multicolumn{2}{c}{MTMD} \\
    \cmidrule(lr){1-2}
    \cmidrule(lr){3-4}
    \cmidrule(lr){5-6}
    \cmidrule(lr){7-8}
    \DAVC{} & \ABM{} & mAP & F1 & mAP & F1 & Acc & F1 \\
    \midrule
     &  & \best{96.0} & 85.2 & \best{75.2} & \best{68.5} & 91.3 & 81.2 \\
     & \cmark & \second{95.9} & \second{85.3} & 73.4 & \second{64.9} & \best{94.3} & \best{87.6} \\
    \cmark &  & 95.8 & 84.5 & 72.8 & 63.9 & \second{93.6} & 84.6 \\
    \cmark & \cmark & 95.8 & \best{86.3} & \second{74.1} & 63.5 & \second{93.6} & \second{86.3} \\
    \bottomrule
  \end{tabular}
  \end{adjustbox}
  % \begin{tablenotes}
  %   \footnotesize
  %   \item[a] MTMD is single-label classification, so the primary metric is accuracy rather than mAP.
  % \end{tablenotes}
  \end{threeparttable}
\end{table}

\section{Discussion}
\label{sec:discussion}

\ToothDINO{} provides a reusable representation that transfers
information from unlabeled panoramic archives to multiple
annotation-limited dental tasks. \DAVC{} and \ABM{} coordinate
anatomical view construction and masked-patch sampling, allowing the
same pretrained encoder to support classification, detection, and
segmentation.

The controlled results show that this strategy is most reliable when
the downstream evidence is concentrated within tooth-bearing anatomy,
whereas gains are smaller or absent for diffuse structures and
heterogeneous image-level labels. The factorial ablation isolates both
\DAVC{} and \ABM{}, but it does not separately decompose the internal
contributions of \DXA{} and \TCC{} within \DAVC{}. The fixed prior may
be less suitable for severe tooth loss, atypical anatomy, or shifted
acquisition geometry. Future work should investigate image-adaptive
priors and independent multi-institutional validation.

\section{Conclusion}
\label{sec:conclusion}

We presented \ToothDINO{}, a geometry-aware pretrained vision encoder obtained by continuing DINOv3 ViT-B/16 on 57,232 unlabeled panoramic radiographs. \DAVC{} and \ABM{} coordinate anatomical view construction and masked-patch sampling during continued pretraining. Across seven datasets and four task families, the matched comparison and five-seed analysis demonstrate stable gains on several downstream tasks, while the smaller improvements on diffuse and heterogeneous targets highlight the task-dependent limitations of the fixed prior. These findings establish \ToothDINO{} as a reusable representation for label-limited panoramic dental image analysis.


\bibliographystyle{IEEEtran}
\bibliography{egbib}

\end{document}
